\documentclass[12pt, a4paper]{article}

% Paquetes esenciales para idioma y formato
\usepackage[utf8]{inputenc}
\usepackage[spanish, es-tabla]{babel}
\usepackage[top=2.5cm, bottom=2.5cm, left=2.5cm, right=2.5cm]{geometry}
\usepackage{hyperref}

% Paquetes para matemáticas y símbolos
\usepackage{amsmath}
\usepackage{amsfonts}

% Paquetes para algoritmos y pseudocódigo
\usepackage{algorithm}
\usepackage{algpseudocode}

% Paquetes para tablas
\usepackage{booktabs}

% Título del Documento
\title{\textbf{Simulación de la propagación de incendios}}
\author{Alejandro Moreno, Facundo Lella, Renzo Dávila, Ramiro Martinez \\ \textit{Programacion Paralela y Distribuida}}
\date{\today}

\begin{document}

\maketitle

\begin{center}
\textbf{Repositorio:} \url{https://github.com/ramirom1/fire-spread-predictor}
\end{center}

\section{Conjunto de datos}

El conjunto de datos que utiliza la simulación se genera proceduralmente al inicio de cada ejecución mediante la función \texttt{createEnvironment(rows, cols, seed)}. La matriz resultante es una grilla bidimensional donde cada celda puede ser Bosque (~70--80\%), Agua (~8--12\%) o Ciudad (~10--15\%), generadas con ruido fractal e interpolación bilineal.

Los experimentos se ejecutan sobre las siguientes combinaciones de tamaño de matriz y cantidad de nodos (procesos MPI), con semilla fija (\texttt{seed = 43}), un máximo de 100\,000 iteraciones, y variando la cantidad de focos iniciales de fuego entre 1, 3 y 5:

\begin{itemize}
    
    \item \textbf{Matriz $10000 \times 10000$} (100 millones de celdas):
    \begin{itemize}
        \item 1 nodo (secuencial)
        \item 2 nodos MPI
        \item 4 nodos MPI
        \item 8 nodos MPI
    \end{itemize}

    \item \textbf{Matriz $15000 \times 15000$} (225 millones de celdas):
    \begin{itemize}
        \item 1 nodo (secuencial)
        \item 2 nodos MPI
        \item 4 nodos MPI
        \item 8 nodos MPI
    \end{itemize}

    \item \textbf{Matriz $20000 \times 20000$} (400 millones de celdas):
    \begin{itemize}
        \item 1 nodo (secuencial)
        \item 2 nodos MPI
        \item 4 nodos MPI
        \item 8 nodos MPI
    \end{itemize}
\end{itemize}

En total se definen $3 \times 4 \times 3 = 36$ configuraciones experimentales. La semilla fija garantiza que, para un mismo tamaño de matriz y cantidad de focos, todos los experimentos operen sobre exactamente el mismo mapa y los mismos puntos de ignición iniciales.

\section{Métricas utilizadas}

Para evaluar el rendimiento se tienen en cuenta las siguientes métricas:

\begin{itemize}
    \item \textbf{Tiempo de ejecución ($T$)}: tiempo de pared medido durante el bucle de simulación. Se usa \texttt{chrono} en la versión secuencial y \texttt{MPI\_Wtime()} en la paralela.

    \item \textbf{Speedup ($S$)}: cuántas veces más rápido corre la versión paralela respecto a la secuencial. $S(p) = T_{\text{sec}} / T_{\text{par}}(p)$. Lo ideal es que sea lineal ($S = p$).

    \item \textbf{Eficiencia ($E$)}: qué tan bien se aprovechan los procesos. $E(p) = S(p) / p$. Un valor cercano a 1 indica buen uso de los recursos.

    \item \textbf{Escalabilidad}: cómo se comporta el sistema al aumentar los procesos o el tamaño del problema. Es importante medirla porque permite saber si tiene sentido agregar más nodos o si el overhead de comunicación empieza a dominar.

    \item \textbf{Balanceo de carga}: en esta simulación el fuego no se distribuye uniformemente en la matriz, por lo que algunos procesos pueden tener mucho más trabajo que otros. Monitorear el desbalance ayuda a entender las pérdidas de eficiencia.

    \item \textbf{Iteración de extinción}: en qué iteración se apaga el fuego. Sirve para verificar correctitud: la versión secuencial y la paralela deben dar el mismo resultado con la misma semilla.
\end{itemize}

\section{Criterios del diseño}

Se varía el tamaño de la matriz ($10000^2$, $15000^2$, $20000^2$), la cantidad de procesos MPI ($1, 2, 4, 8$), y la cantidad de focos iniciales de fuego ($1, 3, 5$), dando un total de 36 combinaciones. La semilla (43) se mantiene fija para que las diferencias de tiempo se deban solo al paralelismo y la cantidad de fuego.

La función de ignición es determinística (basada en semilla, iteración y coordenadas globales), así que la simulación da el mismo resultado sin importar cuántos procesos se usen. Esto permite comparar directamente los tiempos entre configuraciones.

Cada configuración se ejecuta varias veces para tener medidas confiables.

\section{Objetivos}

Los objetivos del experimento son:

\begin{enumerate}
    \item \textbf{Verificar la correctitud de la paralelización}: confirmar que la versión MPI produce resultados idénticos a la versión secuencial para la misma semilla y cantidad de focos, validando que la iteración de extinción coincida en cada tamaño de matriz.

    \item \textbf{Evaluar la escalabilidad fuerte}: para cada tamaño de matriz ($10000^2$, $15000^2$, $20000^2$), medir cómo varía el tiempo de ejecución al incrementar la cantidad de procesos MPI ($p = 1, 2, 4, 8$), y calcular el speedup y la eficiencia en cada caso.

    \item \textbf{Estudiar el efecto del tamaño del problema}: comparar los resultados de speedup y eficiencia entre los tres tamaños de matriz para determinar si problemas más grandes aprovechan mejor la paralelización.

    \item \textbf{Identificar el impacto de las comunicaciones}: entender la importancia y el costo que tienen las comunicaciones entre procesos en el tiempo de ejecución total.

    \item \textbf{Determinar la configuración óptima}: identificar la combinación de tamaño de problema y cantidad de procesos que maximiza el speedup, proporcionando una guía práctica para la ejecución en el clúster.
\end{enumerate}


\end{document}